% ID: SEL-D0015-BA-E02
% ESTADO: cotejada. Segunda pasada visual completa y borrador compilado revisados.
% Fuente: Selectividad/matematicas_aplicadas/matematicas_aplicadas_P_C.pdf
% SHA-256: 463bf63d4697a18a0b9efc9287e7b0a593c803bab98a3cfab0e260067a0dd02c
% Página: 1; bloque A; ejercicio original 2.

\begin{enumerate}[label=\alph*)]
  \item \textbf{(1 punto)} Una fábrica de electrodomésticos dispone de dos
  cadenas de montaje. En una hora de trabajo, la cadena A produce 10 lavadoras
  y 5 frigoríficos, mientras que la cadena B produce 7 lavadoras y 6 frigoríficos.
  El coste de cada hora de trabajo en las cadenas A y B es de $1\,200$ y
  $1\,500$ euros, respectivamente. La cadena A puede funcionar, como máximo,
  el doble de horas que la cadena B. Si deben producir como mínimo 400 lavadoras
  y 280 frigoríficos, formule, \underline{sin resolver}, el problema que permite
  obtener las horas de funcionamiento de las cadenas A y B para minimizar el
  coste de producción de esos electrodomésticos.
  \item \textbf{(1.5 puntos)} Represente el recinto definido por las siguientes
  inecuaciones y calcule sus vértices:
  \[
  \begin{gathered}
  x+2y\geq 7\qquad 4x-y\geq 1\qquad 2x-y\leq 4\\
  3x+2y\leq 20\qquad x\geq 0\qquad y\geq 0
  \end{gathered}
  \]
  Obtenga el valor mínimo de la función $F(x,y)=2x+y$ en el recinto anterior,
  así como el punto en el que se alcanza.
\end{enumerate}
